% Collective memory appendix — drop-in section for songlines_symbolic_memory.tex
%
% Usage:
%   In the main .tex, add at the desired location:
%     % Collective memory appendix — drop-in section for songlines_symbolic_memory.tex
%
% Usage:
%   In the main .tex, add at the desired location:
%     % Collective memory appendix — drop-in section for songlines_symbolic_memory.tex
%
% Usage:
%   In the main .tex, add at the desired location:
%     % Collective memory appendix — drop-in section for songlines_symbolic_memory.tex
%
% Usage:
%   In the main .tex, add at the desired location:
%     \input{collective_memory_appendix.tex}
%
%   This file is self-contained and uses only the styles already loaded
%   by neurips_2026.sty.  It introduces \subsection / \subsubsection
%   but no top-level \section so it can be placed under an existing one.

\section{Collective semantic memory for multi-agent navigation}
\label{sec:collective-memory}

The single-agent symbolic memory of Section~5 is extended to a
multi-agent collective by composing four orthogonal layers.  Each layer
is gated by an independent mode axis, so any subset can be enabled at
test time without retraining or refactoring the base agent.

\subsection{Architecture}
\label{sec:collective-arch}

\paragraph{Phase~1 — event bus.}
A single \emph{collective memory} object stores an append-only log of
typed events.  The dominant event type is
\texttt{place\_observed}$(agent, place\_key, semantic\_tags,
confidence, seq)$.  Recency-weighted retrieval on a tag returns
candidate places ranked by $\lambda^{age}\cdot conf$.

\paragraph{Phase~2 — shared concept graph.}
A \emph{place alignment engine} consolidates events with compatible
spatial and semantic descriptors into shared concepts $c$, each
carrying $(member\_places, dominant\_tag, semantic\_profile,
supporting\_agents)$.  Subsequent queries are answered at the concept
level rather than the raw event level, giving cross-agent generalisation.

\paragraph{Phase~3 — belief dynamics.}
Two parallel mechanisms act on the concept graph between refreshes:
(i) a temporal decay engine multiplies each concept's freshness by
$\rho^{(\text{seq}-\text{last\_seq})}$ with floor $\epsilon$, and
(ii) a conflict rule set scans incompatible tag pairs
$(t^+, t^-)$ and emits per-concept $conflict\_score \in [0,1]$.  The
recall score becomes
$score = conf \cdot tag\_match \cdot \log(1+supp) \cdot freshness \cdot (1-\eta_c)$.

\paragraph{Phase~4 — multi-channel semantic field.}
On top of the canonical graph, a channel-indexed activation field
$A(k,c)$ is maintained per $(channel, concept)$.  The update rule is
\begin{equation}
\label{eq:field-update}
A_{t+1}(k,c) = \lambda\,A_t(k,c)
              + \alpha\,B(c)\,I(k,c)
              + \gamma\,D_t(k,c)
              - \eta\,X_t(c)
              - \xi\,U_t(c),
\end{equation}
where $B(c)$ is a belief strength aggregating confidence/freshness/support/purity;
$I(k,c)$ is a channel affinity from $semantic\_profile$ to channel
weights (with negative weights for competing tags); $D_t$ is a Gaussian
spatial diffusion term; $X_t$ is the $conflict\_score$ from Phase~3;
and $U_t$ is an occupancy/reservation pressure (Section~\ref{sec:phase4c}).
Coefficients $\lambda,\alpha,\eta,\xi,\gamma$ are exposed to the
adaptive layer of Section~\ref{sec:phase4d}.

\subsubsection{Mode axis}

Phase~4 is gated by a single mode parameter taking values
\textsc{none}, \textsc{descriptive}, \textsc{read\_only},
\textsc{coordinated}.  Setting the mode to \textsc{none} reduces the
adapter to Phase~1 raw retrieval, providing a clean baseline.  Modes
upgrade monotonically: \textsc{read\_only} reranks Phase~2/3 candidates
by the field activation; \textsc{coordinated} additionally lets agents
commit reservations that immediately penalise other agents' queries.

\subsection{Reservation-based deconfliction (Phase~4c)}
\label{sec:phase4c}

In \textsc{coordinated} mode an agent may commit a soft reservation
$\langle c, k, agent, ttl \rangle$, which immediately subtracts $\xi$
from $A(k,c)$ and registers in $U_t$.  Subsequent queries from any
agent see the penalised activation without waiting for the next field
refresh.  This is sufficient to deconflict competing claims on the
same resource: when two agents both query for a target tag, the first
to commit a reservation pushes its target's activation below the
runner-up, redirecting the second agent.

\subsection{Outcome-driven adaptive reweighting (Phase~4d)}
\label{sec:phase4d}

A lightweight tracker accumulates per-concept and per-reservation
outcomes over a rolling window.  After every $K$ episodes it applies a
small set of rules:
\begin{itemize}
\item Concept with $base\_conflict \ge 0.15$ and failure rate
      $\ge 0.60$ $\Rightarrow$ $\eta \leftarrow \eta \cdot 1.15$.
\item Reservation success rate $\ge 0.70$ $\Rightarrow$ $\xi \leftarrow \xi \cdot 1.05$;
      rate $< 0.30$ $\Rightarrow$ $\xi \leftarrow \xi \cdot 0.95$.
\item Global failure rate $\ge 0.60$ $\Rightarrow$ $\gamma \leftarrow \gamma \cdot 0.90$.
\end{itemize}
All coefficients are clamped to fixed ranges
($\eta\in[0.10,0.95]$, $\xi\in[0.05,0.60]$, $\gamma\in[0.00,0.30]$).
The update is intentionally non-neural and entirely interpretable.

\subsection{Diagnostic ablations}
\label{sec:collective-ablation}

Table~\ref{tab:collective-smoke} summarises eight unit-level smoke
tests that exercise each phase in isolation, using deterministic
synthetic scenarios so the result is reproducible across runs.

\begin{table}[h]
\centering
\small
\caption{Per-phase diagnostic smokes.  All deterministic; numbers are
from a single run (commit \texttt{a388925}).}
\label{tab:collective-smoke}
\begin{tabular}{lll}
\toprule
Phase & Test & Result \\
\midrule
1     & Recall recovers GT top-1 from event bus & top-1 accuracy = 1.0 \\
2     & Concept consolidation A/B vs. raw events & precision gain $>0$ \\
3a    & Stale observations suppressed by decay & suppression rate = 1.00 \\
3b    & Incompatible tags produce conflict penalty & penalty = 0.74 \\
3c    & Incremental refresh stability & churn = 0.09 \\
4a    & Cross-channel separation, conflict suppression, decay & sep $=0.177$, suppr $=1.0$ \\
4b    & Field reranking on contested target & prec@1 maintained, rank stable \\
4c    & Reservation deconfliction (two consumers, one target) & deconfliction rate = 1.00 \\
4d    & Adaptive reweighting rules fire correctly & $\eta,\xi,\gamma$ adjust within bounds \\
\bottomrule
\end{tabular}
\end{table}

\subsection{End-to-end multi-agent navigation}
\label{sec:collective-multiagent}

To validate the architecture beyond unit-level smokes we built a
two-agent grid environment with two clean water cells located in
opposite corners.  Both agents must reach water; agents block each
other on cell-level collisions.  The same scenario is run under three
field modes, with $N=10$ seeds and a step limit of 80.

\begin{table}[h]
\centering
\small
\caption{Multi-agent navigation under three field modes.
Coordinated mode is the only mode that solves the scenario for both
agents; the failure mode under \textsc{descriptive}/\textsc{read\_only}
is that both agents commit to the same water cell at $t=0$, and the
second agent gets stuck once the first occupies it.}
\label{tab:multiagent-end-to-end}
\begin{tabular}{lcccc}
\toprule
Mode & success (both) & mean steps & dup target rate & hazard hits \\
\midrule
descriptive  & 0.00 & 80.0 & 1.00 & 2.0 \\
read\_only   & 0.00 & 80.0 & 1.00 & 2.0 \\
coordinated  & 1.00 & 21.0 & 0.00 & 4.0 \\
\bottomrule
\end{tabular}
\end{table}

\subsection{Adaptive convergence in episode loop}
\label{sec:collective-adaptive-loop}

The adaptive tracker of Section~\ref{sec:phase4d} is connected to a
20-episode navigation loop with a contested concept whose
$base\_conflict\approx 0.65$.  After every three episodes the tracker
records failures for any high-conflict concept whose channel activation
exceeds a threshold and calls \texttt{adapt()}.
Figure~\ref{fig:eta-trajectory} (text form below) shows the resulting
$\eta_{conflict}$ trajectory.

\begin{verbatim}
episode  0  3  6  9  12 15 18
eta      0.30 0.30 0.345 0.397 0.456 0.525 0.525
\end{verbatim}

The coefficient rises monotonically by $+75\%$ over five effective
adapt cycles and then plateaus once the contested concept's activation
is suppressed below the failure-recording threshold.  The plateau is
the desired convergence behaviour: the loop self-stops once the field
has solved the underlying suppression problem.

\subsection{Invariants}

Four invariants are preserved across all four phases.
\textbf{I1:} the Phase~1 event bus is append-only and is never written
to by higher layers.
\textbf{I2:} the Phase~2 concept graph is the canonical state; the field
of Phase~4 reads from it but never writes.
\textbf{I3:} Phase~3 belief dynamics are applied before each field
refresh; the field's $X_t$ term reads pre-computed
$conflict\_score$.
\textbf{I4:} the planner and local controller of Section~5 are not
modified; the collective layer only reorders candidates and exposes
reservation primitives.  Setting all mode axes to \textsc{none}
recovers the single-agent baseline of Section~5.

These invariants make ablation rigorous: switching any single mode off
is sufficient to attribute an observed effect to a specific phase.

%
%   This file is self-contained and uses only the styles already loaded
%   by neurips_2026.sty.  It introduces \subsection / \subsubsection
%   but no top-level \section so it can be placed under an existing one.

\section{Collective semantic memory for multi-agent navigation}
\label{sec:collective-memory}

The single-agent symbolic memory of Section~5 is extended to a
multi-agent collective by composing four orthogonal layers.  Each layer
is gated by an independent mode axis, so any subset can be enabled at
test time without retraining or refactoring the base agent.

\subsection{Architecture}
\label{sec:collective-arch}

\paragraph{Phase~1 — event bus.}
A single \emph{collective memory} object stores an append-only log of
typed events.  The dominant event type is
\texttt{place\_observed}$(agent, place\_key, semantic\_tags,
confidence, seq)$.  Recency-weighted retrieval on a tag returns
candidate places ranked by $\lambda^{age}\cdot conf$.

\paragraph{Phase~2 — shared concept graph.}
A \emph{place alignment engine} consolidates events with compatible
spatial and semantic descriptors into shared concepts $c$, each
carrying $(member\_places, dominant\_tag, semantic\_profile,
supporting\_agents)$.  Subsequent queries are answered at the concept
level rather than the raw event level, giving cross-agent generalisation.

\paragraph{Phase~3 — belief dynamics.}
Two parallel mechanisms act on the concept graph between refreshes:
(i) a temporal decay engine multiplies each concept's freshness by
$\rho^{(\text{seq}-\text{last\_seq})}$ with floor $\epsilon$, and
(ii) a conflict rule set scans incompatible tag pairs
$(t^+, t^-)$ and emits per-concept $conflict\_score \in [0,1]$.  The
recall score becomes
$score = conf \cdot tag\_match \cdot \log(1+supp) \cdot freshness \cdot (1-\eta_c)$.

\paragraph{Phase~4 — multi-channel semantic field.}
On top of the canonical graph, a channel-indexed activation field
$A(k,c)$ is maintained per $(channel, concept)$.  The update rule is
\begin{equation}
\label{eq:field-update}
A_{t+1}(k,c) = \lambda\,A_t(k,c)
              + \alpha\,B(c)\,I(k,c)
              + \gamma\,D_t(k,c)
              - \eta\,X_t(c)
              - \xi\,U_t(c),
\end{equation}
where $B(c)$ is a belief strength aggregating confidence/freshness/support/purity;
$I(k,c)$ is a channel affinity from $semantic\_profile$ to channel
weights (with negative weights for competing tags); $D_t$ is a Gaussian
spatial diffusion term; $X_t$ is the $conflict\_score$ from Phase~3;
and $U_t$ is an occupancy/reservation pressure (Section~\ref{sec:phase4c}).
Coefficients $\lambda,\alpha,\eta,\xi,\gamma$ are exposed to the
adaptive layer of Section~\ref{sec:phase4d}.

\subsubsection{Mode axis}

Phase~4 is gated by a single mode parameter taking values
\textsc{none}, \textsc{descriptive}, \textsc{read\_only},
\textsc{coordinated}.  Setting the mode to \textsc{none} reduces the
adapter to Phase~1 raw retrieval, providing a clean baseline.  Modes
upgrade monotonically: \textsc{read\_only} reranks Phase~2/3 candidates
by the field activation; \textsc{coordinated} additionally lets agents
commit reservations that immediately penalise other agents' queries.

\subsection{Reservation-based deconfliction (Phase~4c)}
\label{sec:phase4c}

In \textsc{coordinated} mode an agent may commit a soft reservation
$\langle c, k, agent, ttl \rangle$, which immediately subtracts $\xi$
from $A(k,c)$ and registers in $U_t$.  Subsequent queries from any
agent see the penalised activation without waiting for the next field
refresh.  This is sufficient to deconflict competing claims on the
same resource: when two agents both query for a target tag, the first
to commit a reservation pushes its target's activation below the
runner-up, redirecting the second agent.

\subsection{Outcome-driven adaptive reweighting (Phase~4d)}
\label{sec:phase4d}

A lightweight tracker accumulates per-concept and per-reservation
outcomes over a rolling window.  After every $K$ episodes it applies a
small set of rules:
\begin{itemize}
\item Concept with $base\_conflict \ge 0.15$ and failure rate
      $\ge 0.60$ $\Rightarrow$ $\eta \leftarrow \eta \cdot 1.15$.
\item Reservation success rate $\ge 0.70$ $\Rightarrow$ $\xi \leftarrow \xi \cdot 1.05$;
      rate $< 0.30$ $\Rightarrow$ $\xi \leftarrow \xi \cdot 0.95$.
\item Global failure rate $\ge 0.60$ $\Rightarrow$ $\gamma \leftarrow \gamma \cdot 0.90$.
\end{itemize}
All coefficients are clamped to fixed ranges
($\eta\in[0.10,0.95]$, $\xi\in[0.05,0.60]$, $\gamma\in[0.00,0.30]$).
The update is intentionally non-neural and entirely interpretable.

\subsection{Diagnostic ablations}
\label{sec:collective-ablation}

Table~\ref{tab:collective-smoke} summarises eight unit-level smoke
tests that exercise each phase in isolation, using deterministic
synthetic scenarios so the result is reproducible across runs.

\begin{table}[h]
\centering
\small
\caption{Per-phase diagnostic smokes.  All deterministic; numbers are
from a single run (commit \texttt{a388925}).}
\label{tab:collective-smoke}
\begin{tabular}{lll}
\toprule
Phase & Test & Result \\
\midrule
1     & Recall recovers GT top-1 from event bus & top-1 accuracy = 1.0 \\
2     & Concept consolidation A/B vs. raw events & precision gain $>0$ \\
3a    & Stale observations suppressed by decay & suppression rate = 1.00 \\
3b    & Incompatible tags produce conflict penalty & penalty = 0.74 \\
3c    & Incremental refresh stability & churn = 0.09 \\
4a    & Cross-channel separation, conflict suppression, decay & sep $=0.177$, suppr $=1.0$ \\
4b    & Field reranking on contested target & prec@1 maintained, rank stable \\
4c    & Reservation deconfliction (two consumers, one target) & deconfliction rate = 1.00 \\
4d    & Adaptive reweighting rules fire correctly & $\eta,\xi,\gamma$ adjust within bounds \\
\bottomrule
\end{tabular}
\end{table}

\subsection{End-to-end multi-agent navigation}
\label{sec:collective-multiagent}

To validate the architecture beyond unit-level smokes we built a
two-agent grid environment with two clean water cells located in
opposite corners.  Both agents must reach water; agents block each
other on cell-level collisions.  The same scenario is run under three
field modes, with $N=10$ seeds and a step limit of 80.

\begin{table}[h]
\centering
\small
\caption{Multi-agent navigation under three field modes.
Coordinated mode is the only mode that solves the scenario for both
agents; the failure mode under \textsc{descriptive}/\textsc{read\_only}
is that both agents commit to the same water cell at $t=0$, and the
second agent gets stuck once the first occupies it.}
\label{tab:multiagent-end-to-end}
\begin{tabular}{lcccc}
\toprule
Mode & success (both) & mean steps & dup target rate & hazard hits \\
\midrule
descriptive  & 0.00 & 80.0 & 1.00 & 2.0 \\
read\_only   & 0.00 & 80.0 & 1.00 & 2.0 \\
coordinated  & 1.00 & 21.0 & 0.00 & 4.0 \\
\bottomrule
\end{tabular}
\end{table}

\subsection{Adaptive convergence in episode loop}
\label{sec:collective-adaptive-loop}

The adaptive tracker of Section~\ref{sec:phase4d} is connected to a
20-episode navigation loop with a contested concept whose
$base\_conflict\approx 0.65$.  After every three episodes the tracker
records failures for any high-conflict concept whose channel activation
exceeds a threshold and calls \texttt{adapt()}.
Figure~\ref{fig:eta-trajectory} (text form below) shows the resulting
$\eta_{conflict}$ trajectory.

\begin{verbatim}
episode  0  3  6  9  12 15 18
eta      0.30 0.30 0.345 0.397 0.456 0.525 0.525
\end{verbatim}

The coefficient rises monotonically by $+75\%$ over five effective
adapt cycles and then plateaus once the contested concept's activation
is suppressed below the failure-recording threshold.  The plateau is
the desired convergence behaviour: the loop self-stops once the field
has solved the underlying suppression problem.

\subsection{Invariants}

Four invariants are preserved across all four phases.
\textbf{I1:} the Phase~1 event bus is append-only and is never written
to by higher layers.
\textbf{I2:} the Phase~2 concept graph is the canonical state; the field
of Phase~4 reads from it but never writes.
\textbf{I3:} Phase~3 belief dynamics are applied before each field
refresh; the field's $X_t$ term reads pre-computed
$conflict\_score$.
\textbf{I4:} the planner and local controller of Section~5 are not
modified; the collective layer only reorders candidates and exposes
reservation primitives.  Setting all mode axes to \textsc{none}
recovers the single-agent baseline of Section~5.

These invariants make ablation rigorous: switching any single mode off
is sufficient to attribute an observed effect to a specific phase.

%
%   This file is self-contained and uses only the styles already loaded
%   by neurips_2026.sty.  It introduces \subsection / \subsubsection
%   but no top-level \section so it can be placed under an existing one.

\section{Collective semantic memory for multi-agent navigation}
\label{sec:collective-memory}

The single-agent symbolic memory of Section~5 is extended to a
multi-agent collective by composing four orthogonal layers.  Each layer
is gated by an independent mode axis, so any subset can be enabled at
test time without retraining or refactoring the base agent.

\subsection{Architecture}
\label{sec:collective-arch}

\paragraph{Phase~1 — event bus.}
A single \emph{collective memory} object stores an append-only log of
typed events.  The dominant event type is
\texttt{place\_observed}$(agent, place\_key, semantic\_tags,
confidence, seq)$.  Recency-weighted retrieval on a tag returns
candidate places ranked by $\lambda^{age}\cdot conf$.

\paragraph{Phase~2 — shared concept graph.}
A \emph{place alignment engine} consolidates events with compatible
spatial and semantic descriptors into shared concepts $c$, each
carrying $(member\_places, dominant\_tag, semantic\_profile,
supporting\_agents)$.  Subsequent queries are answered at the concept
level rather than the raw event level, giving cross-agent generalisation.

\paragraph{Phase~3 — belief dynamics.}
Two parallel mechanisms act on the concept graph between refreshes:
(i) a temporal decay engine multiplies each concept's freshness by
$\rho^{(\text{seq}-\text{last\_seq})}$ with floor $\epsilon$, and
(ii) a conflict rule set scans incompatible tag pairs
$(t^+, t^-)$ and emits per-concept $conflict\_score \in [0,1]$.  The
recall score becomes
$score = conf \cdot tag\_match \cdot \log(1+supp) \cdot freshness \cdot (1-\eta_c)$.

\paragraph{Phase~4 — multi-channel semantic field.}
On top of the canonical graph, a channel-indexed activation field
$A(k,c)$ is maintained per $(channel, concept)$.  The update rule is
\begin{equation}
\label{eq:field-update}
A_{t+1}(k,c) = \lambda\,A_t(k,c)
              + \alpha\,B(c)\,I(k,c)
              + \gamma\,D_t(k,c)
              - \eta\,X_t(c)
              - \xi\,U_t(c),
\end{equation}
where $B(c)$ is a belief strength aggregating confidence/freshness/support/purity;
$I(k,c)$ is a channel affinity from $semantic\_profile$ to channel
weights (with negative weights for competing tags); $D_t$ is a Gaussian
spatial diffusion term; $X_t$ is the $conflict\_score$ from Phase~3;
and $U_t$ is an occupancy/reservation pressure (Section~\ref{sec:phase4c}).
Coefficients $\lambda,\alpha,\eta,\xi,\gamma$ are exposed to the
adaptive layer of Section~\ref{sec:phase4d}.

\subsubsection{Mode axis}

Phase~4 is gated by a single mode parameter taking values
\textsc{none}, \textsc{descriptive}, \textsc{read\_only},
\textsc{coordinated}.  Setting the mode to \textsc{none} reduces the
adapter to Phase~1 raw retrieval, providing a clean baseline.  Modes
upgrade monotonically: \textsc{read\_only} reranks Phase~2/3 candidates
by the field activation; \textsc{coordinated} additionally lets agents
commit reservations that immediately penalise other agents' queries.

\subsection{Reservation-based deconfliction (Phase~4c)}
\label{sec:phase4c}

In \textsc{coordinated} mode an agent may commit a soft reservation
$\langle c, k, agent, ttl \rangle$, which immediately subtracts $\xi$
from $A(k,c)$ and registers in $U_t$.  Subsequent queries from any
agent see the penalised activation without waiting for the next field
refresh.  This is sufficient to deconflict competing claims on the
same resource: when two agents both query for a target tag, the first
to commit a reservation pushes its target's activation below the
runner-up, redirecting the second agent.

\subsection{Outcome-driven adaptive reweighting (Phase~4d)}
\label{sec:phase4d}

A lightweight tracker accumulates per-concept and per-reservation
outcomes over a rolling window.  After every $K$ episodes it applies a
small set of rules:
\begin{itemize}
\item Concept with $base\_conflict \ge 0.15$ and failure rate
      $\ge 0.60$ $\Rightarrow$ $\eta \leftarrow \eta \cdot 1.15$.
\item Reservation success rate $\ge 0.70$ $\Rightarrow$ $\xi \leftarrow \xi \cdot 1.05$;
      rate $< 0.30$ $\Rightarrow$ $\xi \leftarrow \xi \cdot 0.95$.
\item Global failure rate $\ge 0.60$ $\Rightarrow$ $\gamma \leftarrow \gamma \cdot 0.90$.
\end{itemize}
All coefficients are clamped to fixed ranges
($\eta\in[0.10,0.95]$, $\xi\in[0.05,0.60]$, $\gamma\in[0.00,0.30]$).
The update is intentionally non-neural and entirely interpretable.

\subsection{Diagnostic ablations}
\label{sec:collective-ablation}

Table~\ref{tab:collective-smoke} summarises eight unit-level smoke
tests that exercise each phase in isolation, using deterministic
synthetic scenarios so the result is reproducible across runs.

\begin{table}[h]
\centering
\small
\caption{Per-phase diagnostic smokes.  All deterministic; numbers are
from a single run (commit \texttt{a388925}).}
\label{tab:collective-smoke}
\begin{tabular}{lll}
\toprule
Phase & Test & Result \\
\midrule
1     & Recall recovers GT top-1 from event bus & top-1 accuracy = 1.0 \\
2     & Concept consolidation A/B vs. raw events & precision gain $>0$ \\
3a    & Stale observations suppressed by decay & suppression rate = 1.00 \\
3b    & Incompatible tags produce conflict penalty & penalty = 0.74 \\
3c    & Incremental refresh stability & churn = 0.09 \\
4a    & Cross-channel separation, conflict suppression, decay & sep $=0.177$, suppr $=1.0$ \\
4b    & Field reranking on contested target & prec@1 maintained, rank stable \\
4c    & Reservation deconfliction (two consumers, one target) & deconfliction rate = 1.00 \\
4d    & Adaptive reweighting rules fire correctly & $\eta,\xi,\gamma$ adjust within bounds \\
\bottomrule
\end{tabular}
\end{table}

\subsection{End-to-end multi-agent navigation}
\label{sec:collective-multiagent}

To validate the architecture beyond unit-level smokes we built a
two-agent grid environment with two clean water cells located in
opposite corners.  Both agents must reach water; agents block each
other on cell-level collisions.  The same scenario is run under three
field modes, with $N=10$ seeds and a step limit of 80.

\begin{table}[h]
\centering
\small
\caption{Multi-agent navigation under three field modes.
Coordinated mode is the only mode that solves the scenario for both
agents; the failure mode under \textsc{descriptive}/\textsc{read\_only}
is that both agents commit to the same water cell at $t=0$, and the
second agent gets stuck once the first occupies it.}
\label{tab:multiagent-end-to-end}
\begin{tabular}{lcccc}
\toprule
Mode & success (both) & mean steps & dup target rate & hazard hits \\
\midrule
descriptive  & 0.00 & 80.0 & 1.00 & 2.0 \\
read\_only   & 0.00 & 80.0 & 1.00 & 2.0 \\
coordinated  & 1.00 & 21.0 & 0.00 & 4.0 \\
\bottomrule
\end{tabular}
\end{table}

\subsection{Adaptive convergence in episode loop}
\label{sec:collective-adaptive-loop}

The adaptive tracker of Section~\ref{sec:phase4d} is connected to a
20-episode navigation loop with a contested concept whose
$base\_conflict\approx 0.65$.  After every three episodes the tracker
records failures for any high-conflict concept whose channel activation
exceeds a threshold and calls \texttt{adapt()}.
Figure~\ref{fig:eta-trajectory} (text form below) shows the resulting
$\eta_{conflict}$ trajectory.

\begin{verbatim}
episode  0  3  6  9  12 15 18
eta      0.30 0.30 0.345 0.397 0.456 0.525 0.525
\end{verbatim}

The coefficient rises monotonically by $+75\%$ over five effective
adapt cycles and then plateaus once the contested concept's activation
is suppressed below the failure-recording threshold.  The plateau is
the desired convergence behaviour: the loop self-stops once the field
has solved the underlying suppression problem.

\subsection{Invariants}

Four invariants are preserved across all four phases.
\textbf{I1:} the Phase~1 event bus is append-only and is never written
to by higher layers.
\textbf{I2:} the Phase~2 concept graph is the canonical state; the field
of Phase~4 reads from it but never writes.
\textbf{I3:} Phase~3 belief dynamics are applied before each field
refresh; the field's $X_t$ term reads pre-computed
$conflict\_score$.
\textbf{I4:} the planner and local controller of Section~5 are not
modified; the collective layer only reorders candidates and exposes
reservation primitives.  Setting all mode axes to \textsc{none}
recovers the single-agent baseline of Section~5.

These invariants make ablation rigorous: switching any single mode off
is sufficient to attribute an observed effect to a specific phase.

%
%   This file is self-contained and uses only the styles already loaded
%   by neurips_2026.sty.  It introduces \subsection / \subsubsection
%   but no top-level \section so it can be placed under an existing one.

\section{Collective semantic memory for multi-agent navigation}
\label{sec:collective-memory}

The single-agent symbolic memory of Section~5 is extended to a
multi-agent collective by composing four orthogonal layers.  Each layer
is gated by an independent mode axis, so any subset can be enabled at
test time without retraining or refactoring the base agent.

\subsection{Architecture}
\label{sec:collective-arch}

\paragraph{Phase~1 — event bus.}
A single \emph{collective memory} object stores an append-only log of
typed events.  The dominant event type is
\texttt{place\_observed}$(agent, place\_key, semantic\_tags,
confidence, seq)$.  Recency-weighted retrieval on a tag returns
candidate places ranked by $\lambda^{age}\cdot conf$.

\paragraph{Phase~2 — shared concept graph.}
A \emph{place alignment engine} consolidates events with compatible
spatial and semantic descriptors into shared concepts $c$, each
carrying $(member\_places, dominant\_tag, semantic\_profile,
supporting\_agents)$.  Subsequent queries are answered at the concept
level rather than the raw event level, giving cross-agent generalisation.

\paragraph{Phase~3 — belief dynamics.}
Two parallel mechanisms act on the concept graph between refreshes:
(i) a temporal decay engine multiplies each concept's freshness by
$\rho^{(\text{seq}-\text{last\_seq})}$ with floor $\epsilon$, and
(ii) a conflict rule set scans incompatible tag pairs
$(t^+, t^-)$ and emits per-concept $conflict\_score \in [0,1]$.  The
recall score becomes
$score = conf \cdot tag\_match \cdot \log(1+supp) \cdot freshness \cdot (1-\eta_c)$.

\paragraph{Phase~4 — multi-channel semantic field.}
On top of the canonical graph, a channel-indexed activation field
$A(k,c)$ is maintained per $(channel, concept)$.  The update rule is
\begin{equation}
\label{eq:field-update}
A_{t+1}(k,c) = \lambda\,A_t(k,c)
              + \alpha\,B(c)\,I(k,c)
              + \gamma\,D_t(k,c)
              - \eta\,X_t(c)
              - \xi\,U_t(c),
\end{equation}
where $B(c)$ is a belief strength aggregating confidence/freshness/support/purity;
$I(k,c)$ is a channel affinity from $semantic\_profile$ to channel
weights (with negative weights for competing tags); $D_t$ is a Gaussian
spatial diffusion term; $X_t$ is the $conflict\_score$ from Phase~3;
and $U_t$ is an occupancy/reservation pressure (Section~\ref{sec:phase4c}).
Coefficients $\lambda,\alpha,\eta,\xi,\gamma$ are exposed to the
adaptive layer of Section~\ref{sec:phase4d}.

\subsubsection{Mode axis}

Phase~4 is gated by a single mode parameter taking values
\textsc{none}, \textsc{descriptive}, \textsc{read\_only},
\textsc{coordinated}.  Setting the mode to \textsc{none} reduces the
adapter to Phase~1 raw retrieval, providing a clean baseline.  Modes
upgrade monotonically: \textsc{read\_only} reranks Phase~2/3 candidates
by the field activation; \textsc{coordinated} additionally lets agents
commit reservations that immediately penalise other agents' queries.

\subsection{Reservation-based deconfliction (Phase~4c)}
\label{sec:phase4c}

In \textsc{coordinated} mode an agent may commit a soft reservation
$\langle c, k, agent, ttl \rangle$, which immediately subtracts $\xi$
from $A(k,c)$ and registers in $U_t$.  Subsequent queries from any
agent see the penalised activation without waiting for the next field
refresh.  This is sufficient to deconflict competing claims on the
same resource: when two agents both query for a target tag, the first
to commit a reservation pushes its target's activation below the
runner-up, redirecting the second agent.

\subsection{Outcome-driven adaptive reweighting (Phase~4d)}
\label{sec:phase4d}

A lightweight tracker accumulates per-concept and per-reservation
outcomes over a rolling window.  After every $K$ episodes it applies a
small set of rules:
\begin{itemize}
\item Concept with $base\_conflict \ge 0.15$ and failure rate
      $\ge 0.60$ $\Rightarrow$ $\eta \leftarrow \eta \cdot 1.15$.
\item Reservation success rate $\ge 0.70$ $\Rightarrow$ $\xi \leftarrow \xi \cdot 1.05$;
      rate $< 0.30$ $\Rightarrow$ $\xi \leftarrow \xi \cdot 0.95$.
\item Global failure rate $\ge 0.60$ $\Rightarrow$ $\gamma \leftarrow \gamma \cdot 0.90$.
\end{itemize}
All coefficients are clamped to fixed ranges
($\eta\in[0.10,0.95]$, $\xi\in[0.05,0.60]$, $\gamma\in[0.00,0.30]$).
The update is intentionally non-neural and entirely interpretable.

\subsection{Diagnostic ablations}
\label{sec:collective-ablation}

Table~\ref{tab:collective-smoke} summarises eight unit-level smoke
tests that exercise each phase in isolation, using deterministic
synthetic scenarios so the result is reproducible across runs.

\begin{table}[h]
\centering
\small
\caption{Per-phase diagnostic smokes.  All deterministic; numbers are
from a single run (commit \texttt{a388925}).}
\label{tab:collective-smoke}
\begin{tabular}{lll}
\toprule
Phase & Test & Result \\
\midrule
1     & Recall recovers GT top-1 from event bus & top-1 accuracy = 1.0 \\
2     & Concept consolidation A/B vs. raw events & precision gain $>0$ \\
3a    & Stale observations suppressed by decay & suppression rate = 1.00 \\
3b    & Incompatible tags produce conflict penalty & penalty = 0.74 \\
3c    & Incremental refresh stability & churn = 0.09 \\
4a    & Cross-channel separation, conflict suppression, decay & sep $=0.177$, suppr $=1.0$ \\
4b    & Field reranking on contested target & prec@1 maintained, rank stable \\
4c    & Reservation deconfliction (two consumers, one target) & deconfliction rate = 1.00 \\
4d    & Adaptive reweighting rules fire correctly & $\eta,\xi,\gamma$ adjust within bounds \\
\bottomrule
\end{tabular}
\end{table}

\subsection{End-to-end multi-agent navigation}
\label{sec:collective-multiagent}

To validate the architecture beyond unit-level smokes we built a
two-agent grid environment with two clean water cells located in
opposite corners.  Both agents must reach water; agents block each
other on cell-level collisions.  The same scenario is run under three
field modes, with $N=10$ seeds and a step limit of 80.

\begin{table}[h]
\centering
\small
\caption{Multi-agent navigation under three field modes.
Coordinated mode is the only mode that solves the scenario for both
agents; the failure mode under \textsc{descriptive}/\textsc{read\_only}
is that both agents commit to the same water cell at $t=0$, and the
second agent gets stuck once the first occupies it.}
\label{tab:multiagent-end-to-end}
\begin{tabular}{lcccc}
\toprule
Mode & success (both) & mean steps & dup target rate & hazard hits \\
\midrule
descriptive  & 0.00 & 80.0 & 1.00 & 2.0 \\
read\_only   & 0.00 & 80.0 & 1.00 & 2.0 \\
coordinated  & 1.00 & 21.0 & 0.00 & 4.0 \\
\bottomrule
\end{tabular}
\end{table}

\subsection{Adaptive convergence in episode loop}
\label{sec:collective-adaptive-loop}

The adaptive tracker of Section~\ref{sec:phase4d} is connected to a
20-episode navigation loop with a contested concept whose
$base\_conflict\approx 0.65$.  After every three episodes the tracker
records failures for any high-conflict concept whose channel activation
exceeds a threshold and calls \texttt{adapt()}.
Figure~\ref{fig:eta-trajectory} (text form below) shows the resulting
$\eta_{conflict}$ trajectory.

\begin{verbatim}
episode  0  3  6  9  12 15 18
eta      0.30 0.30 0.345 0.397 0.456 0.525 0.525
\end{verbatim}

The coefficient rises monotonically by $+75\%$ over five effective
adapt cycles and then plateaus once the contested concept's activation
is suppressed below the failure-recording threshold.  The plateau is
the desired convergence behaviour: the loop self-stops once the field
has solved the underlying suppression problem.

\subsection{Invariants}

Four invariants are preserved across all four phases.
\textbf{I1:} the Phase~1 event bus is append-only and is never written
to by higher layers.
\textbf{I2:} the Phase~2 concept graph is the canonical state; the field
of Phase~4 reads from it but never writes.
\textbf{I3:} Phase~3 belief dynamics are applied before each field
refresh; the field's $X_t$ term reads pre-computed
$conflict\_score$.
\textbf{I4:} the planner and local controller of Section~5 are not
modified; the collective layer only reorders candidates and exposes
reservation primitives.  Setting all mode axes to \textsc{none}
recovers the single-agent baseline of Section~5.

These invariants make ablation rigorous: switching any single mode off
is sufficient to attribute an observed effect to a specific phase.
